\documentclass[a4paper,oneside]{book}
\usepackage{fontspec}
\usepackage{bm}
\usepackage[small,sf,bf]{titlesec}
\usepackage{polyglossia}\setdefaultlanguage{english}
\usepackage{amsmath,amssymb}
\usepackage{geometry}
\usepackage[autostyle=true]{csquotes}
\usepackage[backend=bibtex,natbib=true]{biblatex}
\usepackage[
    pdfauthor={Lukas Prokop},
    pdftitle={Bayesian statistics and Neural networks},
    pdfsubject={Bayesian probability theory and Neural networks applied to character recognition},
    pdfkeywords={Bayes, probability, neural network, classifier, character, unicode, kanji, mathematical expressions}
    ]{hyperref} % dep: after biblatex
\usepackage{microtype}
\usepackage{ntheorem}  % dep: before thmtools
\usepackage{thmtools}
\usepackage{xunicode}
\usepackage{xcolor}
\usepackage{metalogo}
\usepackage[fallback]{xeCJK}
\usepackage{mathtools}

\bibliography{thesis}

\newcommand\abs[1]{\left|\hspace{0.6999pt}#1\hspace{0.6999pt}\right|}
\newcommand\set[1]{\left\{#1\right\}}
\newcommand\avg[1]{\overline{#1}}
\newcommand\E{\mathbb E}
\newcommand\mean{\mu}
\newcommand\Nat{\mathbb N}
\newcommand\Prob{\mathbb P}
\newcommand\Var{\mathbb V}
\DeclareMathOperator\Cov{Cov}
\DeclareMathOperator*{\argmin}{arg\,min}
\newcommand\var{\sigma^2}
\newcommand\sd{\sigma}
\newcommand\transpose{T}
\newcommand\card[1]{\left|\hskip .06667em #1\hskip .06667em \right|}

\newcommand\symrepr[1]{\texttt{#1}}
\newcommand\usymrepr[2]{\texttt{#1} ({#2})}
\newcommand\TODO[1]{\textbf{ToDo:} #1 \par}

\renewcommand\theauthor{Lukas Prokop}
\newcommand\authormail{admin@lukas-prokop.at}
\newcommand\institute{Department of Mathematics}
\newcommand\university{University of Kobe, Japan}
\title{Machine Learning using Bayesian statistics and Neural networks}
\author{\theauthor}
\date{Oct 2016 to 2017}

\definecolor{linkcolor}{cmyk}{0.68,0.00,0.71,0.39}
\definecolor{citecolor}{cmyk}{0.60,0.00,0.75,0.51}
\definecolor{urlcolor}{cmyk}{0.59,0.00,0.76,0.44}

\defaultfontfeatures{Ligatures=TeX}%,Numbers=OldStyle}

\setromanfont{Andada}
\setsansfont{DejaVu Sans}
\setCJKmainfont{IPAGothic}
\setCJKsansfont{IPAGothic}
\setCJKfallbackfamilyfont{rm}{IPAGothic}

\newfontfamily\deja{DejaVu Sans}

\hypersetup{
  pdfpagemode={UseOutlines},
  pdftitle={Bayesian statistics and Neural networks},
  pdfauthor={Lukas Prokop},
  pdfkeywords={bachelor thesis, Machine learning, Bayesian statistics, probability theory, Neural networks},
  bookmarksopen=true,
  bookmarksopenlevel=0,
  hypertexnames=false,
  colorlinks=true,
  citecolor=citecolor,
  linkcolor=linkcolor,
  urlcolor=urlcolor,
  pdfstartview={FitV},
  unicode,
  breaklinks=true
}

% theorem design via http://tex.stackexchange.com/a/262573
\newlength\exampleindent
\setlength{\exampleindent}{20pt}
\theorembodyfont{\normalfont}
\theoremindent\exampleindent
\theoremrightindent 10pt
\theoremstyle{plain}

%\newtheoremstyle{mytheoremstyle}% name
%    {\topsep}% Space above
%    {\topsep}% Space below
%    {\itshape}% Body font
%    {}% Indent amount
%    {\normalfont}% Theorem head font
%    {.}% Punctuation after theorem head
%    {.5em}% Space after theorem head
%    {}% Theorem head spec (can be left empty, meaning ‘normal’)

%\theoremstyle{mytheoremstyle}

\newtheorem{example}{Example}
\newtheorem{definition}{Definition}
\newtheorem{theorem}{Theorem}
\newtheorem{proof}{Proof}
\newtheorem{problem}{Problem}

\renewcommand{\baselinestretch}{1.2}

\setlength{\parindent}{0pt}%
\setlength{\parskip}{\baselineskip}%

\begin{document}
\frontmatter
\begin{titlepage}
  \begin{center}
    \textsc{\Large Bachelor Thesis, Mathematics}\\[1cm]

    {\huge \bfseries ``Machine Learning using \\ Bayesian statistics and Neural networks''\\}%
    \vspace{1.7cm}%

    \large \textit{A thesis submitted in fulfillment of the requirements \\ for the bachelor's degree in Mathematics}%
    \vspace{1cm}%

    \begin{minipage}[t]{0.6\textwidth}%
      \centering\large
      \begin{tabular}{rl}
        Author:       & \href{mailto:\authormail}{\theauthor} \\
        Supervisor:   & 高山・信毅 {\deja ※} \\
        Supervisor:   & Bredies Kristian †
      \end{tabular}
    \end{minipage}%
    \vspace{1cm}%

    \begin{minipage}[t]{0.6\textwidth}%
     \centering
      {\deja ※} University of Kobe, Japan \\
      † University of Graz, Austria
    \end{minipage}%
    \vspace{3cm}

    {\large Version: \today}\\[14pt]
    \begin{minipage}{0.4\textwidth}
      \centering
      \includegraphics[width=\textwidth]{img/kobe.pdf}
    \end{minipage}
    \vfill
  \end{center}
\end{titlepage}
\thispagestyle{empty}

\newcommand{\nextoddpage}{
  \if@openright\cleardoublepage\else\clearpage\fi
  \ifdef{\phantomsection}{\phantomsection}{}
}

\newenvironment{declaration}{
  {\nextoddpage}
  %{\space Affidavit}
  {\thispagestyle{plain}}
  {\null\vfil}
  {\noindent%\huge\bfseries{
    \huge{\textsc{Affidavit}}
  %}
  \par\vspace{10pt}}
}{}

\NewDocumentEnvironment{acknowledgements}{}{%
  {\nextoddpage}
  %\tttypeout{\acknowledgementname}
  \thispagestyle{plain}
  \begin{center}{
      %\huge\textit{
      \Large\textsc{Acknowledgements}
      %}
    \par}\end{center}
}{\vfil\vfil\null}

\NewDocumentEnvironment{abstract}{}{%
  {\nextoddpage}
  \thispagestyle{plain}
  \begin{center}{\huge\textsc{Abstract}\par}\end{center}
}{\vfil\vfil\null}

\clearpage
\ifodd\value{page}\else\hbox{}\newpage\fi

\begin{abstract}%
  \parskip5pt
  Machine Learning is a vivid research area. Machine Learning fundamentally changes
  the idea that programmers mechanically write programs in order to perform
  tasks such as classification, regression, clustering, density estimation, and model selection.
  In supervised learning, machines learn by observing test vectors and their desired output. 
  They successively adapt their estimation of the input to return desired outputs for actual data.
  Validation data is used to verify whether this estimate performs good on input,
  the machine has not learned about. Recent efforts such as Google DeepMind's AlphaGo
  or Neural Algorithms of Artistic Style (by Leon A. Gatys, Alexander S. Ecker, and Matthias Bethge)
  show the success of Machine Learning in real-world applications.
  This thesis will cover Bayesian and Neural Networks as two branches of Machine Learning.

  Bayesian networks consider probabilities as quantification of belief.
  Bayes' Theorem is used as a very generic tool to establish a relation between the prior, a likelihood, and the posterior belief.
  For example, consider a set of random variables modelling symptoms of a disease.
  Probabilistic dependencies between these symptoms exist inherently.
  For a set of sample patients, the belief in experiencing certain symptoms and a certain disease is specified.
  The machine can now learn these input-output relationships.
  In the following for a new patient, the belief in symptoms is provided
  and the machine returns a belief whether the patient suffers from a certain disease.

  Neural networks are another branch of Machine Learning. Neural networks consist of multiple
  layers of neurons. Input signals traverse these layers and using sophisticated algorithms,
  neuron weights are adjusted so that neurons learn which input signals belong to which
  output class. This research area is prominent since the 1970s, when Paul Werbos described
  the Backpropagation algorithm in is Ph.D. thesis.

  This bachelor thesis sums up fundamental theorems and theoretical background of the aforementioned fields.
  Furthermore it covers technical details of my implementation to recognize mathematical expressions.

  \vspace{3pt}
  \textbf{Keywords:}
    Machine learning, Bayesian network, Bayesian statistics, Bayes' Theorem, probability theory, Neural networks
\end{abstract}

\begin{abstract}%
  \parskip5pt
  抽象は日本語で仕上がます… \dots

  \vspace{3pt}
  \textbf{Keywords:}
    Machine learning, Bayesian statistics, Bayes' Theorem, Probability theory, Neural networks
\end{abstract}

\begin{acknowledgements}%
  \parskip5pt
  I would like to thank my advisor 高山先生 of Kobe University for his continued effort during my studies.
  Thank you for the valuable input and one year of academic support.

  I would also like to thank my Austrian advisor Bredies Kristian for a final revision and grading
  at University of Graz.

  During my year abroad in Japan, I met lot of new people and gained valuable experiences.
  But I wouldn't have taken the challenge without Martina. Thank you for staying with me
  and sharing your experiences continuously.

  どうもありがとうございました。
\end{acknowledgements}

All source codes are available at \href{http://lukas-prokop.at/proj/bakk\_kobe}{lukas-prokop.at/proj/bakk\_kobe}
and published under terms and conditions of Free/Libre Open Source Software.
This document was printed with \XeLaTeX{} in the Andada typeface.

\tableofcontents
\mainmatter

\chapter{Bayesian theory}
\section{Probability theory and statistics}
\subsection[Sigma-algebra]{$\boldsymbol\sigma$-algebra}
%
The following mathematical object is necessary in order to define probability properly:

\begin{definition}
  Let $X$ be a set. A $\sigma$-algebra is a set $Y$ of subsets of $X$ satisfying:
  \begin{enumerate}
    \item $\emptyset \in Y$
    \item $Z \in Y \implies Z^C \in Y$ where $Z^C$ denotes the complement of $Z$, $X \setminus Z$.
    \item $\left(\bigcup_{i=1}^\infty Z_i\right) \in Y$ where $n \in \Nat$ and $Z_i \in Y$ where $i = 1, 2, \ldots$.
  \end{enumerate}
\end{definition}
When $X$ is a finite set, we may limit the third condition to a finite union.

\begin{example}
  Let $X \coloneqq \set{a, b, c, d}$ and $Z \coloneqq \set{\set{a}}$. We extend $Z$ to a $\sigma$-algebra $Y$:
  \[ Y = \set{\set{}, \set{a,b,c,d}, \set{a}, \set{b,c,d}} \]
\end{example}
The notion of the $\sigma$-algebra is essential to define the notion of the probability space and the random variable rigorously.
In this thesis, our discussion is rigorous when $X$ is a finite set.
When $X$ is an infinite set, we need a full general discussion of measure theory.
Then our discussion will sometimes be intuitive or informal.

\subsection{Basic definitions}
\label{sec:bp-def}
%
Probability theory is concerned with random experiments and random phenomena. Probability in its basic form is the fraction of events with a certain outcome to the total number of events.
%
\begin{definition}
  \label{def:prob}
  \index{Probability space}
  \index{Probability measure}
  \index{Event}
  \index{Random variable}
  A \emph{probability space} $(\Omega, \mathcal A, \Prob)$ denotes
    the set of possible outcomes, a set of events, and a map from an element of $\mathcal A$ to a real value in $[0,1]$.
    $\mathcal A$ is a $\sigma$-algebra. As elements of $\mathcal A$ are sets, we can apply set operations on them.
    Let $\Omega$ be a finite set.
    \emph{Probability measure} $\Prob$ satisfies the following conditions:
    \begin{align}
      &\Prob[A] \geq 0 \text{ for } A \in \mathcal A \label{def:prob-1} \\
      &\Prob[A \cup B] = \Prob[A] + \Prob[B] \text{ for } A, B \in \mathcal A \text{ and } A \cap B = \emptyset \label{def:prob-2} \\
      &\Prob[\Omega] = 1 \label{def:prob-3}
    \end{align}
    These axioms are due to Kolmogorov~\cite{kolmogorov1950foundations}.
    Property~\ref{def:prob-2} implies linearity of the probability measure (for mutually exclusive events $A$):
    \begin{align}
      \Prob\left[\bigcup_{A \in \mathcal A} A\right] &= \sum_{A \in \mathcal A} \Prob\left[A\right] \label{eq:prob-linear}
    \end{align}
  For $R$ \emph{taking the value $k$}, we write,
  \[ \Prob[R=k] \coloneqq \Prob\left(\set{\omega \in \Omega \mid R(\omega) = k}\right) \]
  An \emph{event} is an element $A$ of $\mathcal A$, hence a subset $A$ of $\Omega$.
  A \emph{$\mathbb Z$-valued random variable $R$} is a map from $\Omega$ to $\mathbb Z$
  such that $R^{-1}(z) \in \mathcal A$ for any $z \in \mathbb Z$.
  A \emph{$\mathbb R$-valued random variable $R$} is a map from $\Omega$ to $\mathbb R$
  such that $R^{-1}((r, s]) \in \mathcal A$ for any real numbers $r < s$.
\end{definition}
\begin{example}
  \label{ex:cointoss}
  A coin toss has two possible outcomes, heads (\textit h) or tails (\textit t).
  We consider two coin tosses.
  Then $\Omega = \set{(h, h), (h, t), (t, h), (t, t)}$ and $\mathcal A$ is the powerset of $\Omega$ (i.e. set of all subsets).
  Let $\Prob[A] = \card{A} / \card{\Omega}$, the size of set $A$ divided by $4$.
  Let $R$ be our random variable mapping to $\mathbb Z$ defined as result of the first coin toss (1 represents head, 0 represents tails).
  Then $\Prob[R=1] = \Prob[\setdef{\omega \in \Omega}{R(\omega) = 1}] = \Prob[\set{(h, h), (h, t)}] = \frac24$.
\end{example}

In the following, we will declare random variables, but won't specify the group explicitly.
Either it is obvious from context (because of the numbers we use) or our statements work for both groups.
In that sense, a \emph{random variable} is either {$\mathbb Z$-valued} or {$\mathbb R$}-valued here.

%\subsection{Average Value}
%\label{sec:bp-av}
%%
%\begin{definition}
%  \index{Average value}
%  \index{Population mean}
%  \index{Sample mean}
%  Let $\Omega$ be a finite set and $R$ be a random variable.
%  The \emph{average value} $\avg{R}$ of random variable $R$ is defined as,
%  \begin{align}
%    \avg{R} &\coloneqq \frac{1}{\card{\Omega}} \sum_{a \in \Omega} R(a)
%  \end{align}
%  If all outcomes of the sample space $\Omega$ are considered,
%  we call $\avg{R}$ the \emph{population mean} (denoted $\mu$), otherwise \emph{sample mean}.
%\end{definition}

\subsection{Expected Value}
\label{sec:bp-ev}
%
\begin{definition}
  \index{Expected value}
  \index{Mean value}
  The \emph{expected value} $\E$ (also \emph{mean value}) of a random variable $R$ is defined by:
  \begin{align}
    \E[R]      &\coloneqq \mu
%              = \sum_{a \in \Omega} \Prob[a] \cdot R(\set{a})
%              = \sum_{z \in E} \Prob[R = z] \cdot z \\
               = \sum_{\omega \in \Omega} \Prob[\set{\omega}] \cdot R(\omega) \qquad
    \E_{A}[R]  = \sum_{a \in A} \Prob[\set{a}] \cdot R(a)
  \end{align}
  %where $\Prob[R = z]$ is the probability of the random variable $R$ taking the value $z$.
  %In other words, $\Prob[R = z] \coloneqq \Prob[A_z]$ with $A_z \coloneqq \setdef{e \in \Omega}{R(e) = z}$.
  Let $R$ be a map of $\Omega$ to $E$.
  Applying a function to an expected value means applying the function to the element of $E$.
  Let $f$ be any function with domain $E$, then
  \[
    \E[f(R)]      = \sum_{\omega \in \Omega} \Prob[\set{\omega}] \cdot f(R(\omega)) \qquad
    \E_{A}[f(R)]  = \sum_{a \in A} \Prob[\set{a}] \cdot f(R(a))
  \]
\end{definition}

Let $R$ and $S$ be two random variables. Let $c \in \mathbb R$.
The following properties are satisfied:
\begin{align}
  \E[c]     &= \sum_{\omega \in \Omega} \Prob[\set{\omega}] \cdot c = 1 \cdot c = c \\
  \E[R + c] &= \sum_{\omega \in \Omega} \left(\Prob[\set{\omega}] \cdot (R(\omega) + c)\right) \notag\\
            &= \sum_{\omega \in \Omega} \left(\Prob[\set{\omega}] \cdot R(\omega)\right) + \sum_{\omega \in \Omega} \left(\Prob[\set{\omega}] \cdot c\right) \notag\\
            &= \E[R] + c \cdot \sum_{\omega \in \Omega} \Prob[\set{\omega}] \notag\\
            &= \E[R] + c \cdot 1 = \E[R] + c \\
  \E[R + S] &= \sum_{\omega \in (\Omega_R \cup \Omega_S)}
              \begin{cases}
                \sum_{\omega \in \Omega_R} \Prob[\set{\omega}] \cdot R(\omega) & \text{if } \omega \in \Omega_R \\
                \sum_{\omega \in \Omega_S} \Prob[\set{\omega}] \cdot R(\omega) & \text{if } \omega \in \Omega_S
              \end{cases}\notag\\
            &= \sum_{\omega \in \Omega_R} \Prob[\set{\omega}] \cdot R(\omega) + \sum_{\omega \in \Omega_S} \Prob[\set{\omega}] \cdot R(\omega) \notag\\
            &= \E[R] + \E[S] \label{eq:ERS-RpS} \\
  \E[c \cdot R]
            &= \sum_{\omega \in \Omega} \Prob[\set{\omega}] \cdot (c \cdot R(\omega)) \notag\\
            &= c \cdot \sum_{\omega \in \Omega} \Prob[\set{\omega}] \cdot R(\omega) \notag\\
            &= c \cdot \E[R]
\end{align}
It immediately follows, that this holds true for $\E_A$ for any $A \in \mathcal A$.

\subsection{Equivalence of the continuous probability model}
\label{sec:bp-continuous}
%
\begin{definition}
  \index{Probability Density Function}
  Let $R$ be an $\mathbb R$-valued random variable and
  $f$ be a continuous function defined in $(-\infty, \infty) \subseteq \mathbb R$.
  Let $f$ satisfy the following properties:
  \[
    \Prob[R \leq y] \coloneqq \int\limits_{-\infty}^y f(x) \, dx \hspace{28pt}
    \Prob[z \leq R \leq y] \coloneqq \int\limits_z^y f(x) \, dx \hspace{28pt}
    \Prob[z \leq R] \coloneqq \int\limits_z^\infty f(x) \, dx
  \]
  This establishes a relation between function $f$ and random variable $R$.
  Because $\Prob[R \leq y]$ satisfies properties~\ref{def:prob-1} and \ref{def:prob-3} of probability measures,
  $f$ also satisfies:
  \begin{align}
    f(x) &\geq 0 \qquad \forall x \in (-\infty, \infty) \label{prop:pdf1} \\
       1 &= \int_{-\infty}^{\infty} f(x) \, dx \label{prop:pdf2}
  \end{align}
  $f$ is called a \emph{Probability Density Function} (PDF).
\end{definition}
\begin{example}
  \[ f(x \; | \; \mu, \sigma^2) = \frac{1}{\sqrt{2\pi\sigma^2} } \; e^{ -\frac{(x-\mu)^2}{2\sigma^2} } \]
  The normal distribution function, depending on parameters $\mu$ and $\sigma^2$, is one example for a probability density function.
  The function is introduced in Section~\ref{sec:bp-norm-dist} in detail.
\end{example}

We have seen that the continuous model follows equivalent properties like the discrete model.
This is also true for the expected value.

\begin{definition}
  \index{Expected value}
  The \emph{expected value} of random variable~$R$ %in the continuous model
  is defined as,
  \begin{align}
    \E[R] &\coloneqq \int_{-\infty}^\infty \left(\Prob[R=x] \cdot x\right) \, dx
  \end{align}
  Since intuitively $\int_{-\infty}^\infty \Prob[R=x] \, dx = \int_{\mathbb R} f(x) \, dx$,
  \[ \E[R] = \int_{\mathbb R} x \cdot f(x) \, dx \]
\end{definition}
%
\begin{proof}
  Let $R$ and $S$ be two random variables. Let $c \in \mathbb R$.
  The expected value satisfies:
  %
  \begin{align}
    \E[c]     &\coloneqq c \label{eq:Ec} \\
    \E[R + c] &\coloneqq \int_{\mathbb R} \Prob[R=x] \cdot (x + c) \, dx \notag\\
              &= \int_{\mathbb R} \left(\Prob[R=x] \cdot x + \Prob[R=x] \cdot c\right) \, dx \notag\\
              &= \int_{\mathbb R} \Prob[R=x] \cdot x \, dx + c \cdot \int_{\mathbb R} \Prob[R=x] \, dx \notag\\
              &= \E[R] + c \cdot 1 = \E[R] + c \\
    \E[R + S] &\coloneqq \int_{\mathbb R} \left(\Prob[R=x] \cdot x + \Prob[S=x] \cdot x\right) \, dx \notag\\
              &= \int_{\mathbb R} \Prob[R=x] \cdot x \, dx + \int_{\mathbb R} \Prob[S=x] \cdot x \, dx \notag\\
              &= \E[R] + \E[S] \label{eq:Evv} \\
    \E[c \cdot X] &\coloneqq \int_{\mathbb R} \Prob[R=x] \cdot (x \cdot c) \, dx \notag\\
              &= c \cdot \int_{\mathbb R} \Prob[R=x] \cdot x \, dx
              = c \cdot \E[X] \label{eq:Ecv}
  \end{align}
\end{proof}
Because $\Prob$ and $\E$ provide the same properties in the discrete and continuous case,
we often do not distinguish between these cases. The statements hold true for $\mathbb Z$-valued
as well as $\mathbb R$-valued random variables.

\subsection{Variance and standard deviation}
\label{sec:bp-var-sd}
%
\begin{definition}
  \index{Variance}
  \emph{Variance} quantifies how strong values are spread out from $\E[R]$: %the expected value $\E[R]$:
  \[ \sigma^2 \coloneqq \E\left[(R - \E[R])^2\right] \]
\end{definition}
Considering the entire population, the variance can also quantify over the population mean $\mu$:
\[ \sigma^2 = \Var[R] \coloneqq \E[(R - \mu)^2] \]
In the discrete case, this is equivalent to,
\[ \Var[R] = \E\left[\sum_{\omega \in \Omega} \Prob[\set{\omega}] \cdot \left(R(\omega) - \mu\right)^2\right] \]
and in the continuous case, we have:
\[ \Var[R] = \E\left[\int_{\mathbb R} \Prob[R=x] \cdot \left(x + \mu\right)^2 \, dx\right] \]
\begin{definition}
  \index{Standard deviation}
  The \emph{standard deviation} is defined as its second root:
  \[ \sd = \sqrt{\Var[R]} \]
\end{definition}

\subsection{Covariance}
\label{sec:covar}
%
\index{Covariance}
\emph{Covariance} measures the joint variability of two given random variables.
It is defined as:
\begin{align}
  \Cov[X,Y]
    &= \E[(X - \E[X])(Y - \E[Y])] \label{eq:cov}\\
    &= \E[XY - Y\cdot\E[X] - X\cdot\E[Y] + \E[X]\E[Y]] \notag\\
    &= \E[XY] - \E[Y]\cdot\E[X] - \E[X]\cdot\E[Y] + \E[X]\E[Y] \notag\\
    &= \E[XY] - \E[X]\cdot\E[Y] \label{eq:cov-exy}
\end{align}
%
If $X$ and $Y$ are independent (Section~\ref{sec:bp-ev-indep}), then the covariance is zero.
\begin{align}
  \Cov[X,Y] &= \E[XY] - \E[X] \cdot \E[Y] = \E[X] \cdot \E[Y] - \E[X] \cdot \E[Y] = 0 \label{eq:CovXY0}
\end{align}
We will exploit the following properties:
\begin{align}
  \Cov[X,Y] &= \E[(X - \E[X])(Y - \E[Y])] = \E[(Y - \E[Y])(X - \E[X])] = \Cov[Y,X] \\
  \Cov[X,X] &= \E[(X - \E[X])^2] = \Var[X]
\end{align}
%
Let $X$ be a set of $n$ independent variables $X_{1 \leq i \leq n}$.
Then it holds that:
\begin{align}
  \Var\left[\sum_{i=1}^n X_i\right]
    &= \E\left[\left(\sum_{i=1}^n X_i - \E\left[\sum_{i=1}^n X_i\right]\right)^2\right] \notag \\
    &= \E\left[\left(\sum_{i=1}^n \left(X_i - \E\left[X_i\right]\right)\right)^2\right] \notag \\
    &= \E\left[\sum_{i=1}^n \left(X_i - \E[X_i]\right) \cdot \sum_{j=1}^n \left(X_j - \E[X_j]\right)\right] \notag \\
    &= \E\left[\sum_{j=1}^n \left(\sum_{i=1}^n (X_i - \E[X_i])\right) (X_j - \E[X_j]) \right] \notag \\
    &= \E\left[\sum_{i,j \in [1,n]} \left(X_i - \E[X_i]\right) \left(X_j - \E[X_j]\right)\right] \notag \\
    &= \sum_{i,j \in [1,n]} \E[(X_i - \E[X_i])(X_j - \E[X_j])] \notag \\
    &= \sum_{i,j \in [1,n]} \Cov[X_i, X_j] \notag \\
    &= \sum_{i \in [1,n]} \Cov[X_i, X_i] + \sum_{\substack{i \neq j \\ i,j \in [1,n]}} \Cov[X_i, X_j] \notag \\
    &= \sum_{i=1}^n \Var[X_i] + \sum_{\substack{i \neq j \\ i,j \in [1,n]}} \Cov[X_i, X_j] \notag \\
    &= \sum_{i=1}^n \Var[X_i] + 2 \cdot \sum_{1 \leq i < j \leq n} \Cov[X_i, X_j] \label{eq:sumVarCov}
\end{align}

\subsection{Law of Large Numbers}
\label{sec:bp-lln}
%
The Law of Large Numbers stresses the practical importance of the expected value.

To prove this theorem, we use Chebyshev's Inequality, the Weak Law of Large numbers,
Borel-Cantelli Lemma, and the Strong Law of Large Numbers.
Depending on your requirements of certainty, either the \emph{Weak} Law of Large numbers or
the \emph{Strong} Law of Large Numbers might be considered as Law of Large Numbers.
Our proof structure is based on Craig A. Tracy's~\cite{cnfgen}\footnote{
  Please recognize that there is a small typographical error on page~3.
  \enquote{$S_n(\omega) = 1$ for every $n$} should be \enquote{$X_n(\omega) = 1$ for every $n$}.
}.

\subsubsection{Chebyshev's Inequality}
\label{sec:bp-chebyshev}
%
Consider the continuous case.
Let $R$ be a random variable, $f$ be a PDF over $R$,
$a \in \mathbb R_{\geq 0}$, $p \in \mathbb N$ and let $\E[R^p]$ be defined as follows:
\begin{align*}
  \E[R^p]
          = \int\limits_{\mathbb R} x^p \cdot f(x) \, dx 
          \geq \int\limits_{x \geq a} x^p \cdot f(x) \, dx 
          \geq a^p \int\limits_{x \geq a} f(x) \, dx 
          = a^p \cdot \Prob[R \geq a]
\end{align*}
The discrete case follows immediately.
This concludes the correctness of the following theorem:
\begin{theorem}[Chebyshev's Inequality Theorem]\label{thm:chebyshev}
  Let $R$ be a random variable, $a \in \mathbb R_{\geq 0}$, and $p \in \mathbb N$ is arbitrary.
  Assume $\E[R^p] < \infty$.
  Then it holds that,
  \[ \Prob[R \geq a] \leq \frac{1}{a^p} \E[R^p] \]
\end{theorem}

\subsubsection{Weak Law of Large Numbers}
\label{sec:bp-weak-law}
%
The next theorem is called Weak Law of Large Numbers.
\begin{theorem}[Weak Law of Large Numbers, Bernoulli's Theorem]\label{thm:weak-lln}
  Let $R_i$ be a sequence of independent and identically distributed random variables
  (see section~\ref{sec:bp-iid} for a definition of i.i.d.)
  with common mean $\mean$ and variance $\var$. Let
  \[
      S_n \coloneqq \sum_{i=1}^n R_i \hspace{50pt}
      T_n \coloneqq \frac{S_n}{n} - \mean
  \]
  Then for any $\varepsilon > 0$,
  \[ \lim_{n\to\infty} \Prob[T_n \geq \varepsilon] = 0 \]
\end{theorem}

\begin{proof}
  First, we determine the expected values.
  \begin{align}
      \E[S_n] &= \E\left[\sum_{i=1}^n R_i\right] = \sum_{i=1}^n \E[R_i] = \sum_{i=1}^n \mu = n \cdot \mu & \eqref{eq:Evv} \label{eq:ES_n} \\
%      \E[T_n] &= \E\left[\frac{1}{n}\left(R_1 + R_2 + \ldots + R_n\right) - \mu\right] \notag\\
%              &= \frac1n \left(\E[R_1] + \E[R_2] + \ldots + \E[R_n]\right) - \E[\mu] & \eqref{eq:Ecv} \notag\\
%              &= \frac{n \cdot \mu}{n} - \mu = 0 \label{eq:tn0}
      \E[T_n] &= \E\left[\frac{S_n}{n} - \mu\right] = \frac{\E\left[S_n\right]}{\E\left[n\right]} - \E\left[\mu\right] & \eqref{eq:Evv} \notag\\
              &= \frac{n \cdot \mu}{n} - \mu = 0  & \eqref{eq:ES_n}, \eqref{eq:Ecv} \label{eq:tn0}
  \end{align}
  We also need a result regarding the variance.
  In the continuous and discrete case, it holds that $a^2 \cdot \Var[X] = \Var[a \cdot X]$:
  \begin{align}
    a^2 \cdot \Var[X] &= a^2 \cdot \int (x - \mu)^2 \cdot f(x) \, dx = \int (a \cdot (x - \mu))^2 \cdot f(x) \, dx \notag \\
    a^2 \cdot \Var[X] &= a^2 \cdot \sum_{i=1}^n p_i \cdot (x_i - \mu)^2 = \sum_{i=1}^n p_i \cdot a^2 \cdot (x_i - \mu)^2 = \sum_{i=1}^n p_i \cdot (a \cdot (x_i - \mu))^2 \label{eq:a2V}
  \end{align}
  The relation $\E[X^2] = \Var[X] + \E[X]^2$ holds as well,
  \begin{align}
    \Var[X] &= \E[(X - \mu)^2] = \E[X^2] - \E[2X\mu] + \E[\mu^2] \notag & \eqref{eq:Evv} \\
            &= \E[X^2] - 2 \cdot \E\left[X \cdot \E[X]\right] + \E\left[\E[X]^2\right] & \eqref{eq:Ecv} \notag\\
            &= \E[X^2] - 2 \cdot \E[X]^2 + \E[X]^2 & \eqref{eq:Ecv} \notag\\
            &= \E[X^2] - \E[X]^2 \label{eq:VXisEX2-EX2}
  \end{align}
  We use this result to prove $\Var[T_n] = \frac{\sigma^2}{n}$.
  \begin{align}
    \Var[T_n] &= \E[T_n^2] - \E[T_n]^2 = \E[T_n^2] - 0^2 = \E[T_n^2] & \eqref{eq:VXisEX2-EX2} \eqref{eq:tn0} \label{eq:EVT_n}\\
              &= \E\left[\left(\frac{S_n}{n} - \mu\right)^2\right] = \E\left[\left(\frac{S_n}{n}\right)^2 - 2\frac{S_n}{n}\mu + \mu^2\right] \notag \\
              &= \E\left[\left(\frac{S_n}{n}\right)^2\right] - 2\cdot\E\left[\frac{\mu}{n} S_n\right] + \E\left[\mu^2\right] & \eqref{eq:Evv}\eqref{eq:Ecv} \notag\\
              &= \Var\left[\frac{S_n}{n}\right] + \E\left[\frac{S_n}{n}\right]^2 - 2\mu \cdot \E\left[\frac{S_n}{n}\right] + \mu^2 & \eqref{eq:VXisEX2-EX2}\eqref{eq:Ec} \notag \\
              &= \frac1{n^2} \cdot \Var\left[S_n\right] + \left(\frac1{n} \cdot \E\left[S_n\right]\right)^2 - \frac{2\mu}{n} \cdot (n \cdot \mu) + \mu^2 & \eqref{eq:Ecv}\eqref{eq:a2V} \notag \\
              &= \frac1{n^2} \cdot \left(\Var\left[S_n\right] + \E[S_n]^2\right) - 2\mu^2 + \mu^2 \notag \\
              &= \frac{1}{n^2} \left(\Var\left[\sum_{i=1}^n R_i\right] + \E[S_n]^2\right) - \mu^2 \notag\\
              &= \frac{1}{n^2} \left(\sum_{i=1}^n \Var[R_i] + 2\sum_{\substack{i<j \\ i,j=1}}^n \Cov[R_i,R_j] + \E[S_n]^2\right) - \mu^2 & \eqref{eq:sumVarCov} \notag\\
              &= \frac{1}{n^2} \left(\sum_{i=1}^n \sigma^2 + 0 + (n \cdot \mu)^2\right) - \mu^2 & \eqref{eq:ES_n} & \eqref{eq:CovXY0} \notag\\
              &= \frac{1}{n^2} \left(n \cdot \sigma^2 + n^2 \cdot \mu^2\right) - \mu^2 = \frac{\sigma^2}{n} + \mu^2 - \mu^2 = \frac{\sigma^2}{n}
  \end{align}
\end{proof}
%
%% REMARK: Yes, this is a quote from Wikipedia. Yes, it is true. But it is too short, because we didn't show that the covariance is zero and whether they are uncorrelated or not. So I prefer the approach above
%\index{Bienaym\'e formula}
% This proves the so-called Bienaym\'e formula.
% A shorter approach uses the definition of the mean:
% \[ \Var[\avg{T_n}] = \Var\left(\frac1n \sum_{i=1}^n X_i\right) = \frac1{n^2} \sum_{i=1}^n \Var[X_i] = \frac{\sigma^2}{n} \]
%
%Furthermore, the following equation holds:
%\begin{align}
%    \Var[T_n] &= \E\left[(T_n - \E[T_n])^2\right] = \E\left[(T_n - 0)^2\right] = \E\left[T_n^2\right] \label{eq:EVT_n}
%\end{align}
%
Now we can apply Chebyshev's Inequality Theorem ($R = T_n$, $a = \varepsilon \in \mathbb R, p = 2$):
\begin{align}
  \Prob[\abs{T_n} \geq \varepsilon]
    &\leq \frac{1}{\varepsilon^2} \E\left[\abs{T_n}^2\right]
    = \frac{1}{\varepsilon^2} \Var\left[\abs{T_n}\right]
    = \frac{1}{\varepsilon^2} \frac{\sigma^2}{n}
    & \eqref{eq:EVT_n} \label{eq:cheby-applied}
\end{align}
For any $\varepsilon > 0$ with $n \to \infty$, it holds that
\[ \Prob[\abs{T_n} \geq \varepsilon] \to 0 \]
\[
    \Leftrightarrow \forall \varepsilon > 0:
    \lim_{n\to\infty} \Prob[\abs{T_n} \geq \varepsilon] = 0
\]

This concludes the proof of the Weak Law of Large Numbers.
In order to finish our proof of the Strong Law of Large numbers,
we will use the Borel-Cantelli Lemma, an important result of measure theory.

\subsubsection{Borel-Cantelli Lemma}
\label{sec:bp-borel-cantelli}
%
\begin{lemma}[Borel-Cantelli Lemma]\label{lemma:bcl}
  Let $R_i$ with $1 \leq i < \infty$ be a sequence of events.
  Assume the sum of these probabilities is finite, then it holds that:
  \begin{align}
    \sum_{i=1}^\infty \Prob[R_i] < \infty
    &\implies
    \Prob\left(\limsup\limits_{i\to\infty} R_i\right) = 0
  \end{align}
  So the probability, that the occuring event is an event which occurs infinitely often, is 0.
\end{lemma}

\begin{proof}
  Please consider, that the limsup is defined as,
  \begin{align}
    \limsup\limits_{i\to\infty} R_i &\coloneqq \bigcap_{j=1}^\infty \bigcup_{i \geq j}^\infty R_i \label{eq:limsup-char}
  \end{align}

  The condition requires, that $\sum_{i=1}^\infty \Prob[R_i] < \infty$.
  This statement is equivalent to
  \begin{align}
    \inf_{j \geq 1} \sum_{i=j}^\infty \Prob[R_i] &= 0  \label{eq:inf0}
  \end{align}

  We can make our final conclusion to prove the Borel-Cantelli Lemma:
  \begin{align}
    \Prob\left[\limsup_{i\to\infty} R_i\right]
      &= \Prob\left[\bigcap_{j=1}^\infty \bigcup_{i=j}^\infty R_i\right] & \eqref{eq:limsup-char} \notag\\
      &\leq \inf_{j \geq 1} \Prob\left[\bigcup_{i=j}^\infty R_i\right] & \eqref{eq:prob-linear}\notag\\
      &\leq \inf_{j \geq 1} \sum_{i=j}^\infty \Prob\left[R_i\right] \notag\\
      &= 0 & \eqref{eq:inf0}
  \end{align}
  The result of an intersection of elements creates a set, which
  is an actual subset in any of these sets. However, the infimum
  is not necessarily an element of the set. Hence, an inequality
  is introduced in the second line.
\end{proof}

\subsubsection{Strong Law of Large Numbers}
\label{sec:bp-strong-law}
%
\begin{theorem}[Strong Law of Large Numbers]\label{thm:slln}
  Assume the definitions of Theorem~\ref{thm:weak-lln}.
  Therefore, $R_1, R_2, \ldots$ is an infinite sequence of independent random variables
  with a common distribution ($\mu = \E[R_j]$, $\sigma^2 = \Var[R_j]$). $S_n$ and $T_n$
  are defined. Now consider event $\mathcal E$:
  \[
    \mathcal E = \set{
      \omega \in \Omega: \lim_{n\to\infty} \frac{S_n(\omega)}{n} = \mu
    }
  \]
  Then it holds that
  \[ \Prob[\mathcal E] = 1 \]
\end{theorem}

\begin{proof}
  The following proof assumes $\sigma^2 = \E[R_j^2] < \infty$ and $\E[R_j^4] < \infty$.
  This restriction makes our proof easier, but the less restricted case $\E[R_j] < \infty$ suffices as assumption
  (but this is not proven in this thesis).

  Without loss of generality we assume $\mu = 0$. \\
  If $\mu = 0$ is not satisfied, we consider $P_j \coloneqq R_j - \mu$ instead.

  First of all, we want to give a brief outline of the proof. If it holds that,
  \[ \lim_{n\to\infty} \frac{S_n(\omega)}{n} \neq 0 \]
  then $\exists \varepsilon \in \mathbb R$ with $\varepsilon > 0$ such that for infinitely many $n$
  \[ \frac{S_n(\omega)}{n} > \varepsilon \]
  So to prove the theorem, we will prove that for every $\varepsilon > 0$,
  \[ \Prob[S_n > n \cdot \varepsilon \text{ infinitely often}] = 0 \]
  In the following, this reveals that
  \[ \Prob[\mathcal E] = \Prob\left[\frac{S_n}{n} = 0\right] = 1 \]
  proving Theorem~\ref{thm:slln}. Hence condition $\frac{S_n}{n} = 0$ holds with probability $1$.

  Now we will actually carry out the proof. We define
  \[ A_n = \set{\omega \in \Omega: S_n \geq n \cdot \varepsilon} \]
  and look at $\Prob[A_n]$ using the Chebyshev inequality (Theorem~\ref{thm:chebyshev}) with $p=4$ and $a=(n\cdot\varepsilon)$:
  \[ \Prob[S_n \geq (n\cdot\varepsilon)] \leq \frac{1}{(n\cdot\varepsilon)^4} \E[S_n^4] \]
  We must determine $\E[S_n^4]$ which equals to
  \begin{align*}
    \E\left[\sum_{1 \leq i,j,k,l}^n R_i R_j R_k R_l\right]
      &= \E\left[\sum_{1\leq i}^n \sum_{1\leq j}^n \sum_{1\leq k}^n \sum_{1\leq l}^n R_i R_j R_k R_l\right] \notag\\
%      &= \E\left[(R_1^4 + \ldots + R_1 R_n^3) + (R_2^4 + \ldots) + \ldots + (R_n^4 + \ldots) \right]
      &= \E\left[(R_1^4 + R_1^3 R_2 + \ldots + R_1^4 R_n) + \ldots + (R_n R_1^3 + \ldots + R_n^4)\right]
  \end{align*}
  Here the expected value of an addition of multiplication of random variables is asked for.
  As far as the random variables are i.i.d.,
  we equivalently look for the addition of multiplication of the expected value of a random variable
  (compare with Equation~\ref{eq:ERS-RpS} and Section~\ref{sec:bp-ev-indep}).
  Because $\E[R_i] = 0$, we can remove all terms containing $R_j$ of degree 1 for any $j$.
  These are terms of the structure (assuming $i, j, k$ and $l$ distinct),
  \[
    \E[R_i^3 R_j], \,
    \E[R_i^2 R_j R_k], \,
    \E[R_i R_j R_k R_l]
  \]
  The non-zero terms are $\E[R_i^4]$ and $\E[R_i^2 R_j^2] = \left(\E[R_i^2]\right)^2 = \sigma^4$.
  Now, we need to quantify the occurences of these non-zero terms.
  $\E[R_i^4]$ occurs $n$ times.
  Terms $\E[R_i^2 R_j^2]$ occur $3n \cdot (n-1)$ times,
    as there are $\frac{(n-1) \cdot n}{2}$ ways to choose 2 indices
    and 6 ways to find $R_i^2 R_j^2$.
  In conclusion, we determined,
  \[
    \E\left[S_n^4\right]
      = n \cdot \E\left[R_1^4\right] + 3n \cdot (n - 1) \cdot \sigma^4
      = n \cdot \left( \E\left[R_1^4\right] + 3n \cdot \sigma^4 - 3 \sigma^4 \right)
  \]
  In this expression, $n$ is our only constant occuring with polynomial degree $2$. So for any $n$ sufficiently large, there exists $C \in \mathbb R$ such that
  \begin{align}
    3 \sigma^4 n^2 + \left(\E\left[R_1^4\right] - 3 \sigma^4\right) \cdot n &\leq C \cdot n^2 \label{eq:3sigman2}
  \end{align}
  \[ \Rightarrow \E\left(S_n^4\right) \leq C n^2 \]

  We return back to Chebyshev's inequality
  \[
      \Prob[S_n \geq (n\cdot\varepsilon)]
      \leq \frac{1}{(n\cdot\varepsilon)^4} \E[S_n^4]
      \leq \frac{C \cdot n^2}{\varepsilon^4 \cdot n^2 \cdot n^2}
  \]
  Let $n_0$ be the smallest $n$ such that Equation~\eqref{eq:3sigman2} is satisfied.
  Then it follows that,
  \[
      \sum_{n \geq n_0} \Prob\left[S_n \geq n \cdot \varepsilon\right]
      \leq \sum_{n \geq n_0} \frac{C}{\varepsilon^4 n^2} < \infty
  \]
  With this approach, we skip a finite number of elements of the sum. This does not affect its convergence or divergence.
  Therefore, the conditions to apply the Borel-Cantelli Lemma (Lemma~\ref{lemma:bcl}) are satisfied.
  For every $\varepsilon > 0$ it holds that,
  \[ \Prob\left[S_n \geq n \varepsilon \text{ infinitely often}\right] = 0 \]
\end{proof}

\subsection{Marginal distribution}
\label{sec:bp-marginalization}
%
\begin{definition}
  \index{Marginal probability}
  \index{Marginalization}
  Let $\Omega$ be finite.
  Let $R$ and $S$ be two, discrete random variables. Then we define,
  \[
    \Prob[R=r, S=s] \coloneqq \Prob[\setdef{\omega \in \Omega}{R(\omega) = r \land S(\omega) = s}]
  \]
  This definition enables us to define the \emph{marginal distribution of $R$ and $S$}.
  We also say, we \emph{marginalize $S$ out of $R$}:
  \[ \Prob[R=r] = \sum_{s \in S} \Prob[R=r, S=s] \]
\end{definition}
\begin{example}
  Consider our coin tossing example again.
  Let $R$ be our random variable mapping to $\mathbb Z$ defined as result of the first coin toss
  (1 represents head, 0 represents tails).
  $\mathbb Z$-valued random variable $S$ is defined as result of the second coin toss.
  Then
  \[
    \Prob[R=1, S=0] = \Prob[\setdef{\omega \in \Omega}{R(\omega) = 1 \land S(\omega) = 0}]
                    = \Prob[\set{(h, t)}] = \frac14
  \]
  Marginalization applies, if we query $\Prob[R=1]$:
  \[
    \Prob[R=1] = \Prob[R=1, S=0] + \Prob[R=1, S=1] = \Prob[\set{(h, t)}] + \Prob[\set{(h, h)}] = \frac14 + \frac14 = \frac24
  \]
\end{example}

\subsection{Joint distribution}
\label{sec:bp-joint-distribution}
%
\begin{definition}
  Let $R_i$ be discrete random variables in $1 \leq i \leq n$.
  Joint distribution is given by the \emph{chain rule of probability}:
  \begin{align*}
    \Prob[R_1=r_1, \ldots, R_n=r_n]
        &= \Prob[R_1=r_1] \cdot \Prob[R_2=r_2|R_1=r_1] \cdots \\
        &\cdot \Prob[R_n=r_n|R_1=r_1,R_2=r_2,\ldots,R_{n-1}=r_{n-1}]
  \end{align*}
  If independence (see Section~\ref{sec:bp-indep}) of the random variables is given, it holds that
  \[ \Prob[R_i=r_i|R_1=r_1,\ldots,R_{i-1}=r_{i-1}] = \Prob[R_i=r_i] \qquad \forall i=1,\ldots,n \]
  \[ \Rightarrow \Prob[R_1=r_1, \ldots, R_n=r_n] = \prod_{i=1}^n \Prob[R_i=r_i] \]
\end{definition}
\begin{example}
  In our coin tossing example, we have that:
  \[ \Prob[R=1, S=0] = \Prob[\set{(h, t), (h, h)}] \cdot \Prob[\set{(t, t), (h, t)}] = \frac24 \cdot \frac24 = \frac14 \]
  In the marginalization example, we used the same query and used a different approach to get the same result.
\end{example}

\subsection{Independence}
\label{sec:bp-indep}
%
\index{Independence}
\index{Mutual independence}
\index{Conditional independence}
\begin{definition}
  We assume $\Omega$ is a finite set and consider
  the probability space $(\Omega, \mathcal A, \Prob)$.
  Two $\mathbb Z$-valued random variables $R$ and $S$ are called \emph{(mutually) independent} if
  \[ \Prob[R=r, S=s] = \Prob[R=r] \cdot \Prob[S=s] \]
  for any $r, s \in {\mathbb Z}$ (see, e.g.,~\cite[p.~27]{ito},~\cite[p.~143]{grindstead}).
\end{definition}

Let $A$ be an element of the $\sigma$-algebra ${\mathcal A}$.
We denote by $1_A$ the indicator function of $A$ defined as
\[
  1_A(a) = \begin{cases}
    1 & \text{ if } a \in A \\
    0 & \text{ if } a \not\in A
  \end{cases}
\]
The indicator function $1_A$ defines a ${\mathbb Z}$-valued random variable.
Let $A$ and $B$ be elements of the $\sigma$-algebra ${\mathcal A}$.
The sets (events) $A$ and $B$ are called independent
when the random variables $1_A$ and $1_B$ are independent.

It is important to distinguish between \emph{mutual independence} and \emph{pairwise independence}.
Pairwise independence of events is a weaker condition.

\begin{example}
  Consider our previous coin tossing example (Example~\ref{ex:cointoss}).
  Let us denote coin tosses $ij$, where $i$ and $j$ are the results of the first and second coin toss respectively.
  Define the random variable $R$ as
  \[ R(00)=0, \quad R(01)=0, \quad R(10)=1, \quad R(11)=1 \]
  and define the random variable $S$ (intuitively the result of the second toss) as
  \[ S(00)=0, \quad S(01)=1, \quad S(10)=0, \quad S(11)=1 \]
  Then we have
  \[
    \Prob[R=0, S=0] = 
    \Prob[R=0, S=1] = 
    \Prob[R=1, S=0] = 
    \Prob[R=1, S=1] = 1/4
  \]
  Since $\Prob[R=i] = \frac12$ and $\Prob[S=j] = \frac12$,
  we can see $\Prob[R=i, S=j] = \Prob[R=i] \cdot \Prob[S=j]$.
  Therefore $R$ and $S$ are independent random variables.
  The elements of $\sigma$-algebra
  $A = \set{01, 11} = R^{-1}(1)$ and $B = \set{10, 11} = S^{-1}(1)$
  are independent.
  Intuitively speaking, the event getting $1$ at the first toss
  and the event getting $1$ at the second toss
  are independent.
\end{example}

\subsection{Expected value of independent random variables}
\label{sec:bp-ev-indep}
%
Let $R$ and $S$ be two random variables. We prove $\E[R \cdot S] = \E[R] \cdot \E[S]$
for the discrete and continuous case:
%
\begin{align*}
  \E[R \cdot S] &= \int_{-\infty}^{\infty} \Prob[R=x] \cdot \Prob[S=y] \cdot R(x) \cdot S(y) \, dx \, dy \\
                &= \left(\int_{-\infty}^{\infty} \Prob[R=x] \cdot R(x) \, dx\right) \left(\int_{-\infty}^{\infty} \Prob[S=y] \cdot S(y) \, dy\right) \\
                &= \E[R] \cdot \E[S] \\
  \E[R \cdot S] &= \sum_{a \in \mathcal A_R} \sum_{b \in \mathcal A_S} \left(\Prob[R=a] \cdot \Prob[S=b] \cdot R(a) \cdot S(b)\right) \\
                &= \left(\sum_{a \in \mathcal A_R} \Prob[R=a] \cdot R(a)\right) \cdot \left(\sum_{b \in \mathcal A_S} \Prob[S=b] \cdot S(b)\right) \\
                &= \E[R] \cdot \E[S]
\end{align*}

\subsection{Conditional independence}
\label{sec:bp-cond-indep}
%
Let $R$, $S$, and $T$ be $\mathbb Z$-valued random variables.
The random variables $R$ and $S$ are conditionally independent under the given random variable $T$
if
\[
  \ProbCond{R=r, S=s}{T=t} =
  \ProbCond{R=r}{T=t} \cdot \ProbCond{S=s}{T=t}
\]
for any $r, s, t \in \mathbb Z$ (see, e.g.,~\cite[3.1]{dawid}).
Let
\[
  A_r = R^{-1}(r) \subseteq \Omega \qquad
  B_s = S^{-1}(s) \subseteq \Omega \qquad
  C_t = T^{-1}(t) \subseteq \Omega
\]
Then we can consider $\Prob[A_r]$, $\Prob[B_s]$, and $\Prob[C_t]$ as probability distribution functions associated to $R$, $S$, and $T$ respectively.
We note that we have
\begin{align*}
  \ProbCond{R=r, S=s}{T=t} &= \frac{\Prob[A_r \cap B_s \cap C_t]}{\Prob[C_t]} \\
  \ProbCond{R=r}{T=t} &= \frac{\Prob[A_r \cap C_t]}{\Prob[C_t]} \\
  \ProbCond{S=r}{T=t} &= \frac{\Prob[B_r \cap C_t]}{\Prob[C_t]}
\end{align*}
by the definition of the conditional probability.

\subsection{Bayes' Theorem}
\label{sec:bp-bayes}
%
\begin{theorem}[Bayes' Theorem]
  Let $A$ and $B$ be two events. Let $\Prob[B] \neq 0$. Then:
  \[ \Prob[A|B] = \frac{\Prob[B|A] \cdot \Prob[A]}{\Prob[B]} \]
\end{theorem}
\begin{proof}
  Due to the chain of probability, the following statement holds:
  \[ \Prob[A,B] = \Prob[B|A] \cdot \Prob[A] = \Prob[A|B] \cdot \Prob[B] \]
  Bayes' theorem follows immediately:
  \[ \frac{\Prob[B|A] \cdot \Prob[A]}{\Prob[B]} = \Prob[A|B] \]
\end{proof}

Bayes' Theorem is fundamental to theory we will cover in the following.
$\Prob[A]$ is called \emph{prior probability} and $\Prob[A|B]$ is called \emph{posterior probability} in the Bayesian interpretation.
The names derive from the fact, that $\Prob[A]$ is known beforehand in most applications and $\Prob[A|B]$ is the degree of belief in $A$ after $B$ happened.

\section{Probability distributions}
\label{sec:bp-dist}
%
Probability distributions are templates for probability density functions
satisfying the criteria mentioned in Section~\ref{sec:bp-continuous}.
They are parameterized by one or more variables.
Besides the continuous distributions presented here, discrete distributions also exist.

\subsection{Normal distribution}
\label{sec:bp-norm-dist}
%
\index{Normal distribution}
\index{Gaussian distribution}
\index{Bell curve}
The \emph{Normal distribution} (also \emph{Gaussian distribution} or informally \emph{Bell curve}) is a continuous
probability distribution parameterized by $\mu$ (the mean of the distribution) and $\sigma^2$ (variance).
\[ f(x \,|\, \mu, \sigma^2) = \frac{1}{\sqrt{2 \pi \sigma^2}} e^{-\frac{(x - \mu)^2}{2 \sigma^2}} \]
The distribution is visualized with 4 configurations in Figure~\ref{fig:norm-dist}.

\begin{figure}[p]
  \begin{center}
    \includegraphics[width=0.7\textwidth]{img/normal_distribution.pdf}
    \caption{
      Four different normal distributions with $\mu=0$ and $\sigma^2 = 0.2$ [blue],
      $0.5$ [red], $1.0$ [green], and $5.0$ [yellow]
    }
    \label{fig:norm-dist}
  \end{center}
\end{figure}

\subsection{Beta distribution}
\label{sec:beta-dist}
%
\index{Beta distribution}
The \emph{Beta distribution} is parameterized by $a$ and $b > 0$.
\[ f(x \,|\, a, b) = \frac{1}{B(a, b)} x^{a - 1} (1 - x)^{b-1} \]
where $B(a, b) = \int_0^1 x^{a-1} (1 - x)^{b-1} \, dx$, the so-called beta function.
The distribution is visualized with 4 configurations in Figure~\ref{fig:beta-dist}.

\begin{figure}[p]
  \begin{center}
    \includegraphics[width=0.7\textwidth]{img/beta_distribution.pdf}
    \caption{
      Four different beta distributions:
      $(a=0.5, b=1.5)$ [blue], $(1, 2)$ [red], $(2.5, 4)$ [green], and $(5, 5)$ [yellow]
    }
    \label{fig:beta-dist}
  \end{center}
\end{figure}

\subsection{Independent and identically distributed}
\label{sec:bp-iid}
%
Let $R_i$ be a set of $n$ random variables. $R_i$ are called \emph{independent and
identically distributed} (i.i.d.) if they are mutually independent and each
variable follows the same probability distribution.

\begin{example}
  Three board game players possess two fair dice each.
  Every player is asked to throw them and tell the sum of these two thrown dice.

  All three die sums share the same distribution.
  Sum $2$ is thrown with a chance of $\frac{1}{36}$.
  Sum $7$ is thrown with a chance of $\frac{6}{36}$ (etc.).
  The result of one player does not influence the result of the other player.
  They are mutually independent.
\end{example}

\section{Graphical models}
\label{sec:gm}
%
\index{Bayesian network}
\index{Graphical model}
In order to specify a probabilistic model, you have to provide the following definitions:
\begin{itemize}
  \item A set of random variables and their co-domain ($E$ is either $\mathbb Z$ or $\mathbb R$)
  \item The random variables' dependencies (according to Section~\ref{sec:bp-indep})
  \item You can assign probabilities (conditional probabilities or unconditioned probabilities)
\end{itemize}
%
A graphical model explicitly defines the dependencies using a graph.
Particularly, we looked at \emph{Bayesian networks}, which use a directed acyclic graph (DAG).

Given a Bayesian network, we can make queries (ask for determination of a probability).
Using all the rules about probability theory we learned before, we can derive new knowledge
and might be able to answer the query.

Graphical models furthermore employ two interesting properties.

\subsection{Joint probability of the network}
\label{sec:gm-joint-prob}
%
The joint probability of the Bayesian network with random variables $R_1, R_2, \ldots, R_n$ is given with
\[ \Prob[R_1, \ldots, R_n] = \prod_{i=1}^n \ProbCond{R_i}{\operatorname{ParentNodes}(R_i)} \]
where $\operatorname{ParentNodes}(R_i)$ denotes all the nodes $R_i$ depends on.
This property can be helpful to answer queries.

\subsection[d-separation]{$d$-separation}
\label{sec:gm-d-separation}
%


\subsection{Tool: Bayesian inference}



\TODO{Show symmetry, decomposition, weak union and contraction, via Dawid et al.}

\section{Example: Polynomial curve fitting problem}
\label{sec:bp-curve-fitting}
%
\subsection{The problem}
%
\TODO{visualization}

In the following, we introduce the curve fitting problem similar to Bishop~\cite[p.~4~ff.]{Bishop}.
The problem is defined as follows:

\begin{problem}[Polynomial curve fitting problem]
  Consider a polynomial of arbitrary degree.

  \begin{description}
  \item{Given}
  $x = (x_1, \ldots, x_n)^N$ as a vector of $N$ x-values and
  $t = (t_1, \ldots, t_n)^N$ as the corresponding y-values drawn from the polynomial.
  Furthermore let $E(w)$ be an error function for given polynomial coefficients $w$.

  \item{Find}
  a polynomial with coefficients $w$ which approximates values $t$ minimizing $E(w)$.
  \end{description}
\end{problem}

The degree of the polynomial is purposely unknown.
\emph{Model selection} is a branch of Machine Learning dedicated to finding appropriate models for given problems.
So for polynomial degree choice for our curve fitting problem, we refer to research literature in Model Selection. \TODO{provide useful references for Curve Fitting}
\index{Error function}
An \emph{error function} in this context is any function using $w$, $x$ and $t$ to determine a numeric value representing how much $w$ deviates from $t$ at $x$.
Popular error functions include
\begin{align}
  E(w) &= \frac12 \sum_{n=1}^N \left(y(x_n, w) - t_n\right)^2 \tag{Mean squared error, MSE} \\
  E(w) &= \sqrt{\frac{1}{N} \sum_{n=1}^N (y(x_n, w) - t_n)^2} \tag{Root mean square, RMS} \\
  E(w) &= \frac1N \sum_{n=1}^N (y(x_n, w) - t_n)              \tag{Mean signed deviation, MSD}
\end{align}
% TODO: error function does not need to be derivable, right? But in our design, we derive it to find the gradient to minimize it. So probably it is not required because of backpropagation, but by gradient descent.

\begin{figure}[!h]
  \begin{center}
    \begin{subfigure}[b]{0.45\textwidth}
      \includegraphics[width=\textwidth]{img/interpolation-linear.pdf}
      \caption{Linear interpolation}
      \label{fig:linear-interpolation}
    \end{subfigure}%
    \begin{subfigure}[b]{0.45\textwidth}
      \includegraphics[width=\textwidth]{img/interpolation-nonlinear.pdf}
      \caption{Nonlinear interpolation}
      \label{fig:nonlinear-interpolation}
    \end{subfigure}
    \caption{
      In Figure~\ref{fig:linear-interpolation}, y-values $t_1$ and $t_2$ can be interpolated
      easily using a linear approach. But with more $y$-values (3, in the case of Figure~\ref{fig:nonlinear-interpolation}),
      the problem often requires non-linear solutions.
    }
    \label{fig:interpolation}
  \end{center}
\end{figure}

\subsection{Overfitting}
\label{sec:bp-overfitting}
%
\index{Training data}
\index{Validation data}
\index{Test data}
\index{Overfitting}
Machine Learning distinguishes between a \emph{training} and \emph{validation} dataset as input.
It uses the training set to learn which output is desired for some given input.
Therefore all elements of the training set are labelled such that the error in the output can be quantified.
\emph{Overfitting} describes the situation, when the learning algorithm approximates the output with little error,
but input from the validation set (which contains different inputs) is computed with high error.
So the algorithm perfectly adapted itself to recognize the training data, but performs badly for any other input.

\TODO{visualization}

\subsection{Regularization as countermeasure}

In order to restrict the neural network from overfitting, you penalize it for using
large weights.

\[ E(w) = \frac12 \sum_{n=1}^N (y(x_n, w) - t_n)^2 + \frac{\lambda}{2} \abs{w}^2 \]

Here, the term $\frac{\lambda}{2} \abs{w}^2$ was added where $\abs{w}^2$ denotes the
matrix multiplication
\[ \abs{w}^2 = w^T \cdot w = w_0^2 + w_1^2 + \ldots + w_M^2 \]

% introduce gaussian distribution to every data point

We now model the problem from a probabilistic view:

\subsection{Maximum Likelihood Estimator}
%
% In particular, explain a theoretical reason of adding the term $|w|^2$ for
% the MLE problem in terms of the Baysian theorem.
%
% The Bayesian prior is exponentially to |w|^2. This is a common criticism of the application
% of MLE to Bayesian theory. There is always the prior as assumption, which in this case has computational implications.
% Computational experiments are required to verify the use.
%
The Maximum Likelihood Estimator (MLE) is a technique to estimate the parameters of a probability distribution.
It maximizes the likelihood that the given data actually occurs.

\TODO{Illustrate that the Curve Fitting problem is considered Bayesian here}

\begin{theorem}
  Consider input data $x$, mean $\mu$ and variance $\sigma^2$:
  \[ \ln{\Prob[x | \mu, \sigma^2]} = -\frac1{2\sigma^2} \sum_{n=1}^N (x_n - \mu)^2 - \frac{N}{2} \ln{\sigma^2} - \frac{N}{2} \ln(2\pi) \]
  Then
  $\mu_{\text{ML}} = \frac{1}{N} \cdot \sum_{n=1}^N x_n$ for maximized $\mu$ and \\
    $\sigma_{\text{ML}} = \frac{1}{N}\cdot \sum_{n=1}^N (x_n - \mu_{\text{ML}})^2$ for maximized $\sigma^2$
\end{theorem}

So we want to determine the 2 parameters of a Gaussian distribution, namely $\mu$ and $\sigma^2$, in the maximum likelihood case.
We begin with $\mu$:

\begin{proof}
\begin{enumerate}
  \item Derive $\ln{\Prob[x| \mu, \sigma^2]}$ for $\mu$
    \begin{align*}
      \frac\partial{\partial \mu} \ln{\Prob[x | \mu, \sigma^2]}
      &= \frac\partial{\partial \mu} \left(-\frac{1}{2\sigma^2} \cdot \sum_{n=1}^N (x_n - \mu)^2 - \frac{N}{2} \ln{\sigma^2} - \frac{N}{2} \ln(2\pi)\right) \\
      &= \frac\partial{\partial \mu} \left(-\frac{1}{2\sigma^2} \cdot \sum_{n=1}^N (x_n^2 - 2 x_n \mu + \mu^2) - \frac{N}{2} \ln{\sigma^2} - \frac{N}{2} \ln(2\pi)\right) \\
      &= -\frac{1}{2\sigma^2} \cdot \sum_{n=1}^N (-2x_n + 2\mu) \\
      &= -\frac1{\sigma^2} \cdot \sum_{n=1}^N (\mu - x_n)
    \end{align*}
  \item Set result zero
    \[ 0 = -\frac{1}{\sigma^2} \cdot \sum_{n=1}^N (\mu - x_n) = \sum_{n=1}^N (\mu - x_n) = N \cdot \mu - \sum_{n=1}^N x_n \]
    \[ \implies \mu_{\text{ML}} = \frac1N \cdot \sum_{n=1}^N x_n \qquad \text{commonly called \enquote{sample mean}} \]
\end{enumerate}
\end{proof}

We continue with $\sigma^2$ and use the same approach:

\begin{proof}
  \begin{enumerate}
  \item Derive $\ln{\Prob[x | \mu, \sigma^2]}$ for $\sigma^2$
    \begin{align*}
      \frac{\partial}{\partial \sigma^2} \ln{\Prob[x | \mu, \sigma^2]}
      &= \frac{\partial}{\partial \sigma^2} \left(-\frac{1}{2\sigma^2} \cdot \sum_{n=1}^N (x_n - \mu)^2 - \frac{N}2 \ln{\sigma^2} - \frac{N}2 \ln(2\pi)\right) \\
      &= \frac{1}{2\sigma^4} \cdot \sum_{n=1}^N (x_n - \mu)^2 - \frac{N}{2} \cdot \frac{1}{\sigma^2} \\
      &= \frac{1}{2\sigma^2} \left(\frac{1}{\sigma^2} \cdot \sum_{n=1}^N (x_n - \mu)^2 - N\right)
    \end{align*}
  \item Set result zero
    \begin{align*}
      0 &= \frac{1}{2\sigma^2} \left(\frac{1}{\sigma^2} \cdot \sum_{n=1}^N (x_n - \mu)^2 - N\right) \\
      N \cdot \sigma^2 &= \sum_{n=1}^N (x_n - \mu)^2 \\
      \sigma^2_{\text{ML}} &= \frac{1}{N} \cdot \sum_{n=1}^N (x_n - \mu)^2 \qquad \text{commonly called \enquote{sample variance}}
    \end{align*}
  \end{enumerate}
\end{proof}

And now we derive the precision parameter $\beta$ in the maximum likelihood case:

\begin{theorem}
  Given
  \[ \ln{\Prob[t | x, w, \beta]} = -\frac{\beta}{2} \cdot \sum_{n-1}^N \left(y(x_n,w) - t_n\right)^2 + \frac{N}2 \ln{\beta} - \frac{N}{2} \ln(2\pi) \]
  then find
  \[ \frac{1}{\beta_{\text{ML}}} = \frac{1}{N} \cdot \sum_{n=1}^N (y(x_n, w_{\text{ML}}) - t_n)^2 \] by maximizing $\beta$
\end{theorem}

\begin{proof}
\begin{enumerate}
  \item Derive $\ln{\Prob[t | x,w,\beta]}$ with $\beta$
    \begin{align*}
      \frac{\partial}{\partial \beta} \ln{\Prob[t | x,w,\beta]}
      &= \frac{\partial}{\partial \beta} \left(-\frac{\beta}{2} \sum_{n=1}^N (y(x_n,w) - t_n)^2 + \frac{N}{2} \ln\beta - \frac{N}{2} \ln(2\pi)\right) \\
      &= -\frac12 \cdot \sum_{n=1}^N \left(y(x_n,w) - t_n\right)^2 + \frac{N}{2} \cdot \frac1\beta
    \end{align*}
  \item Set result zero
    \begin{align*}
      0 &= -\frac{1}{2} \cdot \sum_{n=1}^N \left(y(x_n,w) - t_n\right)^2 + \frac{N}{2\beta} \\
      \frac{N}{\beta} &= \sum_{n=1}^N \left(y(x_n,w) - t_n\right)^2 \\
      \frac{1}{\beta_{\text{ML}}} &= \frac{1}{N} \cdot \sum_{n=1}^N (y(x_n,w) - t_n)^2
    \end{align*}
\end{enumerate}
\end{proof}







The maximum of the logarithm of an expression corresponds to the minimum of the negative logarithm of the same expression. \TODO{so why do we minimize and not maximize?}

\begin{align}
  -\log{\Prob[\omega|x,t,\alpha,\beta]}
    &\propto -\log\left[\Prob[t|x,\omega,\beta] \cdot \Prob[\omega|\alpha]\right] \\
\intertext{due to proportionality, $\exists c \in \mathbb R$ such that} \\
    &= -\log{\Prob[t|x,\omega,\beta]} - \log\Prob[\omega|\alpha] - \log{c} \\
\intertext{insert formula Bishop 1.62} \\
    &= \frac\beta2 \sum_{n=1}^N \left(y(x_n,\omega)-t_n\right)^2 - \frac{N}{2} \ln\beta + \frac{N}{2} \ln{2\pi} - \log\left(\left(\frac{\alpha}{2\pi}\right)^{\frac{M+1}{2}} \cdot \exp\left(-\frac{\alpha}{2} \omega^\transpose \omega\right)\right) - \log{c} \\
    &= \frac\beta2 \sum_{n=1}^N \left(y(x_n,\omega)-t_n\right)^2 - \frac{N}{2} \ln\beta + \frac{N}{2} \ln{2\pi} - \frac{M+1}{2} \log\alpha + \frac{M+1}{2} \log{2\pi} + \frac{\alpha}{2} \omega^\transpose \omega - \log{c}
\end{align}

Let $f$ be any function with a minimum. Then $\argmin_{\omega} f(\omega) = \argmin_{\omega} c \cdot f(\omega) + a$ for any $c, a \in \mathbb R$.
This applies also to our case:

\begin{align}
  \argmin_{\omega} -\log{\Prob[\omega|x,t,\alpha,\beta]}
    &= \argmin_{\omega} \left(\frac\beta2 \cdot \sum_{n=1}^N \left(y(x_n, \omega) - t_n\right)^2 + \frac\alpha2 \omega^\transpose \omega\right) \\
    &= \argmin_{\omega} \beta\left(\frac12 \sum_{n=1}^N \left(y(x_n, \omega) - t_n\right)^2 + \frac{\frac\alpha\beta}{2} \omega^\transpose \omega^\transpose\right) \\
    &= \argmin_{\omega} \frac12 \sum_{n=1}^N \left(y(x_n, \omega) - t_n\right)^2 + \frac{\frac\alpha\beta}{2} \omega^\transpose \omega^\transpose \\
\end{align}

Hence, the coefficients maximizing the probability that the coefficients correspond to our model parameters ($x, t, \alpha, \beta$) are given in the last line.
Considering we determine the best coefficients by minimizing the error function, it is justified to consider these coefficients as optimum. Let $\lambda = \frac{\alpha}{\beta}$, then \dots

\[ \tilde{E}(\omega) = \frac12 \sum_{n=1}^N \left(y(x_n, \omega) - t_n^2\right)^2 + \frac{\lambda}{2} \|\omega\|^2 \]


\chapter{Neural networks}

%\input{appendix}

\backmatter
\printbibliography
%\listoffigures
%\listoftables
%\printindex

\end{document}
